% GENERATED FILE. Edit canonical lessons, then run npm run generate:books.
% canonical-source-hash: fnv1a64:b36db1c92a08b6d6
% canonical-lessons: TA-C126-eri, TA-C126-kulam, TA-C126-katu, TA-C126-ilai, TA-C126-man

\chapter{Lake, Pond, Forest, Leaf, Soil}
\label{ch:ta-lake-pond-forest-leaf-soil}
\hlchaptermodality{126}

\begin{chapteropening}
\textbf{By the end of this chapter:} \emph{I can say a lake, a pond, a forest, a leaf and soil.}

Complete the last lesson of chapter 126: i can say a lake, a pond, a forest, a leaf and soil.
\end{chapteropening}

\section[ēri]{\ta{ஏரி} (ēri) — a lake}

\label{lesson:TA-C126-eri}



\begin{quote}
Before the new one: say the Tamil for coffee, then the Tamil for the sea.
\end{quote}

\subsection*{You'll want to know: \ta{ஏரி}}
\textbf{\ta{ஏரி}} — \emph{ēri} — \textquotedblleft{}a lake\textquotedblright{}.

\begin{grammarlens}[title={What you've built}]
One more word for this chapter.
\end{grammarlens}

\subsection*{Your turn}
\begin{itemize}
\raggedright
  \item \emph{Say it:} \emph{ēri}
  \item \emph{Say it:} \emph{ēri}, once more
  \item \emph{Say it:} say \emph{ēri}
  \item \emph{Recall:} say the Tamil for coffee, then the Tamil for the sea, then say \emph{ēri} again
\end{itemize}

\begin{tcolorbox}[breakable,colback=teal!4,colframe=teal!35!black,title={Before you move on}]
What does \ta{ஏரி} mean? (\textquotedblleft{}A lake\textquotedblright{}.) Say it once more.
\end{tcolorbox}
\section[kuḷam]{\ta{குளம்} (kuḷam) — a pond}

\label{lesson:TA-C126-kulam}



\begin{quote}
Before the new one: say the Tamil for the sea, then the Tamil for a lake.
\end{quote}

\subsection*{You'll want to know: \ta{குளம்}}
\textbf{\ta{குளம்}} — \emph{kuḷam} — \textquotedblleft{}a pond\textquotedblright{}.

\begin{grammarlens}[title={What you've built}]
One more word for this chapter.
\end{grammarlens}

\subsection*{Your turn}
\begin{itemize}
\raggedright
  \item \emph{Say it:} \emph{kuḷam}
  \item \emph{Say it:} \emph{kuḷam}, once more
  \item \emph{Say it:} say \emph{kuḷam}
  \item \emph{Recall:} say the Tamil for the sea, then the Tamil for a lake, then say \emph{kuḷam} again
\end{itemize}

\begin{tcolorbox}[breakable,colback=teal!4,colframe=teal!35!black,title={Before you move on}]
What does \ta{குளம்} mean? (\textquotedblleft{}A pond\textquotedblright{}.) Say it once more.
\end{tcolorbox}
\section[kāṭu]{\ta{காடு} (kāṭu) — a forest}

\label{lesson:TA-C126-katu}



\begin{quote}
Before the new one: say the Tamil for a lake, then the Tamil for a pond.
\end{quote}

\subsection*{You'll want to know: \ta{காடு}}
\textbf{\ta{காடு}} — \emph{kāṭu} — \textquotedblleft{}a forest\textquotedblright{}.

\begin{grammarlens}[title={What you've built}]
One more word for this chapter.
\end{grammarlens}

\subsection*{Your turn}
\begin{itemize}
\raggedright
  \item \emph{Say it:} \emph{kāṭu}
  \item \emph{Say it:} \emph{kāṭu}, once more
  \item \emph{Say it:} say \emph{kāṭu}
  \item \emph{Recall:} say the Tamil for a lake, then the Tamil for a pond, then say \emph{kāṭu} again
\end{itemize}

\begin{tcolorbox}[breakable,colback=teal!4,colframe=teal!35!black,title={Before you move on}]
What does \ta{காடு} mean? (\textquotedblleft{}A forest\textquotedblright{}.) Say it once more.
\end{tcolorbox}
\section[ilai]{\ta{இலை} (ilai) — a leaf}

\label{lesson:TA-C126-ilai}



\begin{quote}
Before the new one: say the Tamil for a pond, then the Tamil for a forest.
\end{quote}

\subsection*{You'll want to know: \ta{இலை}}
\textbf{\ta{இலை}} — \emph{ilai} — \textquotedblleft{}a leaf\textquotedblright{}.

\begin{grammarlens}[title={What you've built}]
One more word for this chapter.
\end{grammarlens}

\subsection*{Your turn}
\begin{itemize}
\raggedright
  \item \emph{Say it:} \emph{ilai}
  \item \emph{Say it:} \emph{ilai}, once more
  \item \emph{Say it:} say \emph{ilai}
  \item \emph{Recall:} say the Tamil for a pond, then the Tamil for a forest, then say \emph{ilai} again
\end{itemize}

\begin{tcolorbox}[breakable,colback=teal!4,colframe=teal!35!black,title={Before you move on}]
What does \ta{இலை} mean? (\textquotedblleft{}A leaf\textquotedblright{}.) Say it once more.
\end{tcolorbox}
\section[maṇ]{\ta{மண்} (maṇ) — soil, earth}

\label{lesson:TA-C126-man}



\begin{quote}
Before the new one: say the Tamil for a forest, then the Tamil for a leaf.
\end{quote}

\subsection*{You'll want to know: \ta{மண்}}
\textbf{\ta{மண்}} — \emph{maṇ} — \textquotedblleft{}soil, earth\textquotedblright{}.

\begin{grammarlens}[title={What you've built}]
One more word for this chapter.
\end{grammarlens}

\subsection*{Your turn}
\begin{itemize}
\raggedright
  \item \emph{Say it:} \emph{maṇ}
  \item \emph{Say it:} \emph{maṇ}, once more
  \item \emph{Say it:} say \emph{maṇ}
  \item \emph{Recall:} say the Tamil for a forest, then the Tamil for a leaf, then say \emph{maṇ} again
\end{itemize}

\begin{tcolorbox}[breakable,colback=teal!4,colframe=teal!35!black,title={Before you move on}]
What does \ta{மண்} mean? (\textquotedblleft{}Soil, earth\textquotedblright{}.) Say it once more.
\end{tcolorbox}
